\chapter{Riot API, adatforrások és adatgyűjtési stratégia}

A rendszer elsődleges adatforrása a Riot Developer Portal és a kapcsolódó League of Legends API-készlet \cite{riotportal,riotapi}. Ezek lehetővé teszik játékosfiókok azonosítását, meccslisták lekérdezését, részletes meccs-JSON-ok letöltését, és résztvevői statisztikák feldolgozását.

Az API ugyanakkor nem korlátlan: a fejlesztői kulcsok szigorú sebességkorlátokkal működnek, és ezek az architektúrát is befolyásolják. Ez magyarázza, hogy a projektben miért különül el világosan az adatgyűjtő, a normalizáló, a gráfépítő és a futtatókörnyezeti réteg. A cél nem az, hogy minden egyszerre, ugyanabban a kérésben történjen, hanem az, hogy a rendszer különálló, újrafuttatható lépésekből álljon össze.

\section{A használt Riot API-végpontok logikája}

A jelenlegi projekt a Riot API több rétegét használja, de a legfontosabb tervezési döntés az volt, hogy a rendszer központi azonosítójává a \textit{puuid} váljon. Ez a döntés végigvonul az egész kódbázison. A felhasználóbarát név és tag ugyan hasznos a frontendben, de az adatbázis és a gráfépítés szempontjából az azonosíthatóság elsődleges.

A pipeline logikailag az alábbi lépésekre bontható:

\begin{enumerate}
    \item egy kiinduló játékos vagy seed-lista alapján a rendszer fiókinformációt kér le;
    \item a fiókból meccsazonosító-listát szerez;
    \item a meccsazonosítókhoz letölti a részletes match JSON-okat;
    \item a résztvevői sorokból \textit{puuid}-alapú játékosprofilt és kapcsolati adatot képez;
    \item az eredményt datasethez kötve perzisztálja.
\end{enumerate}

Ez azért szerencsés, mert a névváltozások, régiós sajátosságok és meccslistás lekérdezések nem keverednek össze a későbbi gráf- és pontszámítási lépésekkel. A \textit{puuid} egyszerre stabil összekötő kulcs a nyers meccsadat, a játékosprofil és a gráfcsúcs között.

\subsection{Régiós routing és végpont-topológia}

A collector az \textit{europe.api.riotgames.com} regionális útvonalon hívja a \textit{match-v5} végpontcsaládot. Két végpont adja a gyűjtés alaplépéseit: a meccsazonosítók lekérése játékosonként, majd az azonosítókhoz tartozó teljes meccs-JSON letöltése. Ez első látásra technikai apróságnak tűnhet, de a teljes adatgyűjtési megbízhatóság múlik rajta: hibás routing mellett a rendszer úgy tűnhet működőképesnek, hogy közben a lekérdezések egy része konzisztensen elbukik. Az adatok lokálissá tételéig a pipeline nem végez elemzést; a feldolgozás teljes egészében lokális fájlokon és adatbázison zajlik.

\subsection{Collector állapotgép és premade-probe}

A \textit{backend/match\_collector.py} nem egyszerű lineáris letöltőszkript, hanem paraméterezhető állapotgép. A konfiguráció tartalmaz dataset-azonosítót, útvonalakat, gyűjtési módot, queue-típust, meccsszám-korlátot, iterációs limitet és premade-küszöböt. Ennek következtében a gyűjtés explicit módon szabályozható, újrafuttatható és különböző célokra hangolható.

A tipikus futás így írható le:

\begin{enumerate}
    \item a backend kijelöli az aktív datasetet és előkészíti az útvonalakat;
    \item a collector véletlen vagy konkrét \textit{puuid} alapján jelölt játékost választ;
    \item rövid meccslista-lekéréssel premade-probe történik;
    \item ha a jelölt elég ismétlődő csapattársi mintázatot mutat, teljesebb gyűjtés indul;
    \item a letöltött meccsek eseményes logolás mellett fájlba kerülnek;
    \item az eredmény később normalizáló és gráfépítő pipeline-ba lép tovább.
\end{enumerate}

Ez a részlet azért fontos, mert a premade-fókusz nem csak a gráfelemzés szintjén jelenik meg, hanem már a mintavételi mechanizmusban is. A rendszer tehát nem utólag próbálja kiszűrni a visszatérő kapcsolatok jelentőségét, hanem a gyűjtési stratégiát is erre hangolja.

\begin{figure}[H]
\centering
\includegraphics[width=0.98\textwidth,height=0.74\textheight,keepaspectratio]{collector-state-machine.pdf}
\caption{A collector paraméterezhető állapotgépének folyamatábrája}
\end{figure}

\section{Gyűjtési fókusz és queue-választás}

A projekt két kutatási datasetet használ: a \textit{flexset} Ranked Flex Queue meccsekből áll, a \textit{soloq} Ranked Solo/Duo Queue meccsekből. A korai fejlesztési fázisban keletkezett \textit{default} dataset vegyes összetételű (Ranked Flex, Ranked Solo/Duo, ARAM, normál draft és egyéb queue-k keverednek benne), ezért kutatási célra nem alkalmas. Az aktív elemzések kizárólag a \textit{flexset} és \textit{soloq} datasetekre támaszkodnak, amelyek célzott és reprodukálható stratégiával épültek fel.

A queue-típus megválasztása nem tetszőleges. A Ranked Flex megengedi a premade csapatok részvételét, ezért a visszatérő szövetségesi kapcsolatok természetes módon jelenhetnek meg. A Ranked Solo/Duo szándékosan kizárja az 5-fős premade-et, így kontrollált összehasonlítási alapként szolgál. Ez a kettősség az empirikus fejezetek értelmezési keretét is megadja.

\section{Premade-fókusz és gyűjtési stratégiai intuíció}

A projekt mögött végig ott állt az a gyakorlati intuíció, hogy az igazán érdekes kapcsolatok nem az egyszeri, véletlenszerű találkozásokból születnek, hanem a visszatérő együtt- vagy egymás ellen játszásból. Emiatt a gyűjtőszkriptek és a dokumentáció egy része kifejezetten premade-szerű vagy ismétlődő kapcsolatok felderítésére van hangolva.

Ez a szemlélet a teljes későbbi rendszerre hatással volt:

\begin{itemize}
    \item a gráfépítésben élküszöb jelent meg;
    \item az analitikák az ismétlődő kapcsolatokat tekintik komoly jelöltnek;
    \item a pathfinder-költségek is a kapcsolaterősséget veszik figyelembe.
\end{itemize}

Másképp fogalmazva: a gyűjtési stratégia nem előszoba, hanem a gráfjelentés egyik meghatározó forrása.

Ez a döntés ugyanakkor tudatos torzítást is bevezet. Ha a gyűjtés a visszatérő, premade-szerű kapcsolatokat részesíti előnyben, akkor a korpusz kevésbé reprezentálja az egyszeri és lazább találkozásokat. A jelen dolgozat céljai mellett ez elfogadható kompromisszum, mert a kutatási kérdés középpontjában éppen a tartósabb kapcsolati magok állnak. A háló tehát nem a teljes populáció semleges lenyomata, hanem az ismétlődő interakciók felé célzottan dúsított részvilág.

\section{Rate limit, megbízhatóság és operatív kompromisszumok}

A Riot API használata az architektúrában nem pusztán külső függőség, hanem korlátozó tényező is. A fejlesztői kulcsok lekérdezésszámhoz kötöttek, ezért a gyűjtőrétegben explicit sebességkorlát-kezelés, várakozási logika és részfeladatokra bontás jelent meg. A dokumentációban rögzített rate-limit elemzés ezt nem adminisztratív részletként, hanem rendszertervezési kényszerként kezeli.

Ennek közvetlen következményei:

\begin{itemize}
    \item a meccsletöltés nem kérésenkénti frontend-akcióként, hanem háttérfolyamatként fut;
    \item a nyers adatállapotot érdemes lokálisan cache-elni;
    \item a grafikus és analitikai oldal nem függhet attól, hogy egy élő API-lekérés éppen sikeres-e.
\end{itemize}

Ez a döntés később az egész rendszer egyik alapelve lett: \emph{a demo- és kutatási réteg a már összegyűjtött korpuszra épül, nem online API-kérésekre}.

Az implementáció szintjén ez több konkrét védelmet jelent. A collector saját másodpercenkénti várakozást alkalmaz, külön kétperces ablakban számlálja a kéréseket, HTTP 429 esetén harminc másodpercet vár, majd újrapróbálkozik, és a már elmentett meccseket fájllétezés alapján kihagyja. Ez a működés idempotenssé teszi a letöltést: egy megszakadt vagy újraindított futás nem írja felül feleslegesen a korábbi állományokat, és nem épít újabb bizonytalanságot a nyers korpuszba.

\begin{table}[H]
\centering
\caption{A collector fő megbízhatósági mechanizmusai}
\begin{tabularx}{\textwidth}{>{\raggedright\arraybackslash}p{4.4cm}|X}
\toprule
\textbf{Mechanizmus} & \textbf{Funkció} \\
\midrule
Másodpercenkénti késleltetés & A rövid távú API-terhelés kontrollja \\
Kétperces számláló & Illeszkedés a Riot kulcshasználati limitjeihez \\
429 újrapróbálás & Robusztusabb működés átmeneti limitütközéskor \\
Létező meccsfájl kihagyása & Duplikált mentés és újra-letöltés elkerülése \\
Eseményes stdout log & Backend- és frontend-oldali állapotkövetés támogatása \\
\bottomrule
\end{tabularx}
\end{table}

Ezek a döntések nem puszta üzemeltetési kényelmi funkciók. A reprodukálható számítógépes kutatás egyik alapelve, hogy az eredményekhez vezető lépések nyomon követhetők és lehetőleg automatizálhatók legyenek \cite{sandve2013}. A collector parametrizálhatósága, eseményes logolása és idempotens mentési logikája pontosan ezt a célt szolgálja.

\section{A flexset és a soloq gyűjtési stratégiájának összehasonlítása}

A \textit{flexset} dataset kizárólag 440-es queue ID-jű (Ranked Flex) meccseket tartalmaz, EUNE Master$+$ ranglista-pozíciójú seed-játékosoktól. A collector premade-probe mechanizmusa kiszűri azokat a jelölteket, akiknek nincs elegendő visszatérő csapattársi mintázatuk; a befogadottakat csapattársaik azonosítói alapján bővíti tovább (snowball sampling). Az eredmény közel 4980 meccs, erősen összefüggő, visszatérő kapcsolatokban gazdag háló.

A \textit{soloq} dataset 420-as queue ID-jű (Ranked Solo/Duo) meccsekből áll, EUNE Master$+$ Solo/Duo ranglistás seed-populációval. Ez nem azonos a \textit{flexset} seed-populációjával: a Flex Queue és a Solo/Duo Queue ranglistái eltérő mechanikán alapulnak, és más játékosbázist fednek le. A probe-logika itt nem szűr premade-jelenlét alapján, a corpus körülbelül 2000 meccsel zárul. A \textit{soloq} szerepe kontrollalap: olyan dataset, ahol a visszatérő premade-kapcsolatok szisztematikusan nem favorizáltak, ezért a két dataset összehasonlítása megalapozottabb.

\begin{table}[H]
\centering
\caption{A két dataset gyűjtési stratégiájának összehasonlítása}
\begin{tabularx}{\textwidth}{>{\raggedright\arraybackslash}p{4.2cm}|X|X}
\toprule
\textbf{Dimenzió} & \textbf{flexset} & \textbf{soloq} \\
\midrule
Queue ID & 440 (Flex) & 420 (Solo/Duo) \\
Premade-probe & Igen, ismétlések alapján & Nem \\
Seed-terjeszkedés & Csapattársak alapján & Listalekéréssel \\
Gyűjtési cél & Szociális struktúra & Teljesítmény-alapvonal \\
Corpus mérete & 4980 meccs & 2000 meccs \\
Szerep a projektben & Fő kutatási adatkészlet & Kontroll-adatkészlet \\
\bottomrule
\end{tabularx}
\end{table}

Mindkét dataset gyűjtése leáll, ha elérte a konfigurált meccsszám- vagy iterációs limitet, a gyűjtési sor kiürült, vagy manuálisan le lett állítva. A nyers meccs-JSON fájlok fájlrendszeren perzisztálódnak; az idempotens fájlellenőrzés garantálja, hogy már letöltött meccsek nem kerülnek újra letöltésre.

\section{Adatminőség, torzítás és óvatos interpretáció}

A Riot API formálisan pontos, de a dolgozat szempontjából az adat így sem „semleges". A mintavételt erősen befolyásolja, hogy milyen seed játékosokról indult a gyűjtés, mely queue-k kerültek előtérbe, milyen mélyre ment a meccslista-lekérés, és a premade-kapcsolatok milyen minimumismétlés után tekinthetők érdekesnek.

Ezért a dolgozat minden későbbi eredményénél fontos megjegyzés, hogy a hálózat \emph{nem a teljes League of Legends univerzum reprezentációja}, hanem egy adott gyűjtési stratégia által létrehozott, vizsgálható részvilág. Ez nem gyengeség, hanem korrekt keretfeltétel.

További torzításforrás, hogy a dataset fokozatosan épülő archívumként működik. A meccsek nem egyetlen, szigorúan lezárt időablakból származnak, ezért a játékosok \textit{match\_count} értékei részben azt is tükrözik, mennyire került rájuk fókusz a gyűjtési folyamatban. Emiatt a score-réteg és az időbeliség kapcsolatát mindig óvatosan kell kezelni: a magasabb mintaszám nem kizárólag aktivitást, hanem gyűjtési figyelmet is jelenthet.
